% ==============================================================================
% CHAPITRE 3 -- DÉMARCHE MÉTHODOLOGIQUE QUALITATIVE
% Mémoire MAE -- Arthur SAUNIER -- Wavestone Luxembourg
% Le \chapter{} est déclaré dans le fichier maître ; ce fichier démarre en \section.
% ==============================================================================

\section{Justification du design de recherche}

La question centrale de ce mémoire interroge la manière dont l'automatisation des
tâches d'analyse redéfinit la légitimité et le développement des compétences des
consultants. Sa formulation en « comment » n'appelle ni la mesure d'un écart ni la
vérification d'une relation entre variables~: elle appelle la compréhension d'un
processus en cours, saisi dans les pratiques et dans les représentations de ceux qui
le vivent. C'est cette nature de la question qui commande le choix méthodologique, et
non l'inverse.

Une démarche quantitative supposerait que les dimensions à mesurer soient déjà
stabilisées. Or elles ne le sont pas. Construire un questionnaire reviendrait à
décider par avance ce qu'est un usage, ce qu'est une perte de compétence ou ce qu'est
une atteinte à la légitimité, puis à demander aux répondants de se positionner sur des
catégories dont rien ne garantit qu'elles recouvrent leur expérience. L'objet est
émergent~; les outils concernés se sont diffusés dans les pratiques plus vite que le
vocabulaire permettant de les décrire. Une enquête par questionnaire produirait donc
une mesure précise de catégories mal fondées.

Deux obstacles pratiques renforcent ce choix. Le premier tient à la taille de la
population~: l'effectif de la practice concernée ne permet aucune exploitation
statistique digne de ce nom, la puissance d'un test étant hors d'atteinte sur quelques
dizaines d'individus. Le second tient à la sensibilité du sujet. Les usages
informels, le recours à des outils non validés par le cabinet ou la validation d'un
contenu que l'on n'aurait pas su produire soi-même ne se déclarent pas dans une case à
cocher. Ces pratiques ne s'obtiennent que par le récit, dans un cadre où le répondant
comprend que la description d'un contournement n'est pas un aveu.

La posture épistémologique retenue est interprétativiste. Elle admet que la réalité
étudiée n'est pas donnée indépendamment des acteurs qui la construisent, et que le
matériau accessible au chercheur est le sens que ces acteurs attribuent à leurs propres
pratiques \cite{gavardperret2018}. Le discours recueilli n'est donc pas traité comme
un reflet transparent de l'activité réelle, mais comme une reconstruction située, dont
les silences, les hésitations et les contradictions internes font partie du matériau
plutôt qu'ils n'en altèrent la qualité.

Le raisonnement suivi est abductif. Le cadrage théorique du chapitre précédent fournit
les catégories de départ~: modes de conversion des connaissances, registres de
légitimité, mécanismes de coordination, déterminants de l'acceptation. Ces catégories
orientent la construction de l'instrument de collecte sans prétendre épuiser ce que le
terrain donnera à voir. L'allée et venue entre les concepts mobilisés et le matériau
recueilli constitue le mouvement analytique propre à ce travail.

L'entretien semi-directif individuel a été préféré aux autres techniques qualitatives
disponibles. Le groupe de discussion aurait exposé les propos les plus sensibles à la
censure des pairs et à la présence hiérarchique. L'observation directe aurait mal
saisi des usages qui se déroulent pour l'essentiel dans un rapport solitaire à un
écran. L'entretien semi-directif, en revanche, articule une liberté de formulation
suffisante pour laisser émerger l'inattendu et une structure commune sans laquelle la
comparaison entre profils serait impossible \cite{blanchet2007}.

Il importe enfin d'énoncer ce que ce design autorise et ce qu'il interdit. Il autorise
l'identification de mécanismes, la mise au jour de tensions entre niveaux
hiérarchiques et la formulation d'hypothèses transférables à des organisations
comparables. Il interdit toute généralisation statistique~: aucun résultat ne pourra
être présenté comme valant pour la profession du conseil, ni même pour l'ensemble du
cabinet. Il interdit également toute conclusion causale et toute mesure de
performance. Si un répondant décrit un gain de temps, ce gain de temps est un fait de
discours et non un fait mesuré~; le dispositif établit ce qui est perçu, vécu et
raconté, non ce qui est produit. Cette frontière est rappelée à chaque étape de la
restitution.

\section{Dispositif empirique et échantillonnage}

L'échantillon visé compte six à huit collaborateurs de Wavestone Luxembourg, tous
rattachés à la practice Cybersécurité. Il ne procède d'aucun tirage aléatoire mais
d'un échantillonnage raisonné, construit par contraste de profils. La variable
discriminante retenue est la position hiérarchique, parce qu'elle détermine le rapport
même à l'objet étudié~: le consultant junior éprouve sa propre montée en compétences,
le manager l'observe chez ceux qu'il encadre, le directeur en arbitre les conséquences
sur le modèle de production du cabinet. Un même phénomène est ainsi saisi depuis trois
positions d'énonciation distinctes.

\begin{table}[h]
\centering
\small
\begin{tabular}{|p{0.13\linewidth}|p{0.09\linewidth}|p{0.32\linewidth}|p{0.32\linewidth}|}
\hline
\textbf{Profil} & \textbf{Codes} & \textbf{Position dans le dispositif} & \textbf{Apport attendu} \\
\hline
Consultants juniors & J1--J3 & Producteurs directs des livrables d'analyse & Usage vécu, sentiment sur l'apprentissage, rapport aux règles \\
\hline
Seniors et managers & M1--M3 & Encadrants de production et responsables du contrôle qualité & Observation de l'évolution des juniors, rôle en revue, relation client directe \\
\hline
Directeurs & D1--D2 & Responsables du modèle économique et de la politique d'outillage & Arbitrage risque-productivité, proposition de valeur, staffing et formation \\
\hline
\end{tabular}
\caption{Composition du panel et logique de contraste}
\label{tab:panel-chap3}
\end{table}

Le tableau~\ref{tab:panel-chap3} explicite la logique de construction du panel. Trois
juniors, trois seniors ou managers et un à deux directeurs permettent de constituer
deux groupes comparables en effectif, condition d'une lecture croisée, tout en
réservant une position de surplomb. Le tiraillement attendu entre le groupe des
juniors et celui des managers est une propriété voulue du design et non un résultat~:
il conditionne la possibilité même d'observer un écart, mais ne préjuge en rien de son
existence ni de son sens. Des critères secondaires --- ancienneté dans le cabinet,
diversité des types de missions --- guident le choix des individus à l'intérieur de
chaque strate, afin d'éviter que le panel ne se réduise à un même segment d'activité.

Le guide d'entretien, reproduit intégralement en annexe~A, est construit de manière
déductive à partir des quatre questions directrices. Chaque thème constitue la
traduction conversationnelle d'une de ces questions. Un thème de cadrage préalable,
consacré à la trajectoire du répondant et à la nature de ses missions, ouvre
l'entretien~: il installe la parole, recueille les variables de contrôle nécessaires à
la lecture croisée et permet à l'enquêteur de repérer le vocabulaire propre au
répondant, qu'il réemploiera ensuite. Le thème~1 porte sur l'usage réel des outils
dans le travail d'analyse, le thème~2 sur l'apprentissage du métier et la construction
de l'expertise, le thème~3 sur la relation au client et la valeur de l'expertise, le
thème~4 sur l'encadrement, le contrôle qualité et la gestion des risques. Un thème de
clôture, projectif et ouvert, invite le répondant à formuler une décision qu'il
prendrait à la place de la direction et à signaler ce que l'entretien n'a pas abordé.

Deux partis pris commandent la formulation des questions \cite{blanchet2007}. Chaque
thème s'ouvre par un récit ou un incident critique plutôt que par une question
d'opinion, le récit d'une situation vécue produisant un matériau vérifiable là où
l'opinion appelle la réponse de façade. L'enquêteur s'interdit par ailleurs de
prononcer le vocabulaire de la recherche --- légitimité, montée en compétences,
hallucination, gouvernance --- avant que le répondant ne l'ait lui-même mobilisé, tout
emploi prématuré induisant la réponse et détruisant la valeur probante du verbatim. La
structure du guide demeure identique pour tous les répondants, condition de la
comparabilité~; seules les formulations d'ancrage sont adaptées au profil, afin de
solliciter la position d'énonciation propre à chacun.

L'entretien est calibré sur trente-cinq à quarante-cinq minutes, durée compatible avec
les contraintes de disponibilité de consultants en mission. Les entretiens se tiennent
en visioconférence ou en présentiel, mais toujours au moyen d'une réunion organisée
dans l'environnement collaboratif du cabinet~: l'enregistrement et la transcription
automatique sont assurés par cet outil, ce qui garantit que l'audio ne quitte jamais
le système d'information de l'organisation. Une seule langue est retenue par entretien,
l'alternance linguistique dégradant fortement la qualité de la transcription
automatique.

\section{Modalités de collecte et d'analyse}

Chaque entretien s'ouvre par la lecture d'un préambule reproduit en annexe~A. Ce
préambule précise l'objet du travail, rappelle qu'il ne s'agit ni d'un audit ni d'une
évaluation, énonce les modalités d'anonymisation et sollicite explicitement
l'autorisation d'enregistrer. En cas de refus, l'entretien se poursuit sous prise de
notes manuscrite et la restitution des propos est alors signalée comme reconstituée.
Les répondants sont désignés par un code de profil suivi d'un rang~; les noms de
personnes, de clients et de missions sont remplacés par des désignations génériques
entre crochets dès la retranscription. Une fiche de traçabilité est renseignée
immédiatement après chaque entretien, avant toute retranscription, afin de conserver
les éléments contextuels que l'enregistrement ne capte pas~: climat de l'échange,
réticences perçues, thèmes écourtés, hypothèses émergentes à tester lors des entretiens
suivants.

La transcription automatique produite par l'outil de visioconférence fait l'objet d'une
reprise manuelle intégrale. Cette relecture n'est pas une simple correction
orthographique~: elle rétablit les hésitations, les reprises et les ruptures de
construction, qui sont autant d'indices analytiques. Le corpus final est constitué des
transcriptions ainsi corrigées, à l'exclusion de toute reformulation.

Le traitement repose sur une analyse thématique de contenu \cite{bardin2001}. Le
matériau est découpé, catégorisé, puis interprété selon les trois temps classiques de
la préanalyse, de l'exploitation et de l'inférence. L'unité de codage retenue est le
segment de sens, c'est-à-dire l'extrait de discours porteur d'une idée complète,
indépendamment de sa longueur grammaticale. Un même segment peut recevoir plusieurs
codes lorsqu'il articule plusieurs dimensions, un consultant décrivant par exemple dans
la même phrase un gain de productivité et une inquiétude sur sa propre compétence.

La grille de codage, reproduite en annexe~B, est mixte. Son ossature est déductive~:
quatre axes correspondent terme à terme aux quatre questions directrices et à leurs
ancrages théoriques respectifs. L'axe~A traite de l'appropriation des outils dans le
travail d'analyse, en distinguant les tâches déléguées, les tâches sanctuarisées, les
canaux d'apprentissage de l'outil et les déterminants d'usage. L'axe~B traite de
l'apprentissage, des compétences et de la transmission, en distinguant la valeur
formatrice attribuée aux tâches, les effets déclarés sur le raisonnement et les
transformations de la relation entre junior et senior. L'axe~C traite de la légitimité
et de la relation client, en distinguant les registres de confiance invoqués, le
traitement de la transparence vis-à-vis du client et les déplacements de la
proposition de valeur. L'axe~D traite de la gouvernance, du contrôle qualité et des
risques, en distinguant le statut des règles, les dispositifs de vérification, les
incidents rencontrés, l'attribution des responsabilités et la perception du cadre
lui-même, jugé tantôt bridant, tantôt insuffisamment explicite. Un cinquième axe,
transversal, est partiellement inductif~: il accueille les écarts entre discours et
pratique, les divergences entre niveaux hiérarchiques, les projections sur l'avenir du
métier, les recommandations spontanées et une catégorie ouverte destinée aux thèmes
non anticipés. Ces derniers sont intitulés et définis en cours d'analyse.

Le codage est conduit au fil de l'eau, chaque entretien étant traité avant la
conduite du suivant. Cette contrainte n'est pas d'ordre pratique~: elle est la
condition d'un suivi honnête de la saturation théorique \cite{glaser1967}. La grille
est stabilisée après le troisième entretien, puis appliquée rétroactivement à
l'ensemble du corpus afin que tous les entretiens soient codés selon le même
référentiel. La saturation est considérée comme atteinte lorsqu'un entretien
supplémentaire n'ajoute aucun code nouveau et ne modifie pas la structure de la
grille~; le rang de l'entretien à partir duquel aucun code inédit n'apparaît est
mentionné dans la restitution. Les contradictions internes à un même entretien ne sont
pas arbitrées mais codées comme telles~: elles constituent un résultat.

Le corpus codé alimente une matrice de restitution croisant les codes en lignes et les
répondants en colonnes, complétée d'une colonne de fréquence et d'un verbatim
représentatif \cite{miles2003}. La lecture en colonnes restitue la cohérence interne
du discours de chaque répondant~; la lecture en lignes fait apparaître les
convergences du panel et ses lignes de fracture, une colonne d'écart signalant les
codes dont la fréquence diffère sensiblement entre le groupe des juniors et celui des
managers. Cette matrice constitue le support direct du chapitre suivant~: chaque
verbatim y cité est référencé par le code du répondant et le code thématique, de
manière à ce que le chemin analytique entre le corpus brut et l'interprétation proposée
demeure reconstituable par le lecteur.

\section{Limites et biais du dispositif}

La taille de l'échantillon constitue la première limite. Six à huit entretiens
autorisent la profondeur, non la représentativité. Rien ne garantit d'ailleurs que la
saturation soit effectivement atteinte au terme du panel prévu~; si des codes inédits
continuent d'apparaître au dernier entretien, ce fait sera signalé comme tel plutôt que
dissimulé, et les conclusions correspondantes présentées comme provisoires.

Le caractère mono-site constitue la deuxième limite. L'étude porte sur une practice
d'une seule entité, sur un marché dont la structure --- place financière concentrée,
pression réglementaire spécifique --- n'est pas transposable telle quelle. Les résultats
sont indexés à ce contexte. Leur portée relève de la transférabilité analytique vers
des organisations comparables, non de la généralisation \cite{lincoln1985}.

La désirabilité sociale constitue la troisième limite, et elle est ici plus vive
qu'ailleurs. Le sujet touche à des pratiques qu'un collaborateur n'a pas intérêt à
exposer~: le recours à des outils hors du cadre officiel, le contournement d'une règle,
l'aveu d'une dépendance ou la validation d'un contenu que l'on n'aurait pas su produire.
Plusieurs dispositions visent à contenir ce biais~: le préambule dissocie explicitement
l'entretien de toute démarche d'évaluation, l'anonymisation est décrite avant
l'enregistrement, les questions procèdent par récit d'incident plutôt que par
sollicitation d'opinion, et l'écart entre position déclarée et pratique décrite est
codé plutôt que résolu. Ces précautions atténuent le biais sans le supprimer.

La position de l'enquêteur constitue la limite la plus lourde et appelle un examen
séparé. Celui-ci est lui-même consultant du cabinet étudié~: il occupe simultanément la
place de l'observateur et celle du participant \cite{girin1989}. Ce double statut n'est
pas un accident
du dispositif, il en est une condition d'existence, et il produit des effets dans les
deux sens.

Les avantages sont réels. L'accès au terrain est immédiat, là où un chercheur externe
se heurterait aux règles de confidentialité qui régissent l'activité de conseil. La
familiarité avec le vocabulaire professionnel, avec le déroulement d'une mission et
avec la chaîne de production d'un livrable permet de comprendre l'implicite, de relancer
au bon endroit et de repérer une réponse de façade. La confiance préexistante entre
collègues rend possible l'évocation de pratiques qu'un inconnu n'obtiendrait pas.

Les risques sont symétriques. Le répondant s'adresse à un collègue et ajuste son
propos en conséquence, particulièrement lorsque l'écart hiérarchique joue en sens
inverse de la situation d'entretien, un directeur restant en position d'autorité face à
un enquêteur qui lui est subordonné. L'enquêteur, de son côté, aborde le terrain avec
des prénotions constituées par sa propre expérience et risque de sélectionner ce qui
les confirme. La connivence professionnelle produit enfin un effet insidieux~: la
formule « tu vois ce que je veux dire » clôt l'explication au moment précis où elle
devient nécessaire, et l'enquêteur qui voit effectivement ce que le répondant veut dire
perd le verbatim qui l'aurait établi.

Quatre mesures visent à contenir ces risques. Le statut de l'enquêteur est explicité
en préambule, de sorte que l'ambiguïté de la situation soit posée plutôt que subie.
L'explicitation de l'implicite est systématiquement demandée, y compris --- et surtout ---
lorsque l'enquêteur a compris~: toute allusion est reprise par une demande de récit
factuel. Une neutralité axiologique stricte est tenue, l'enquêteur s'abstenant de
défendre comme de critiquer l'usage de ces outils, toute marque d'approbation orientant
durablement la suite de l'entretien. La frontière entre les deux statuts est enfin
maintenue dans le traitement~: seul le matériau enregistré est codé, la connaissance de
terrain servant à interpréter les propos mais jamais à établir un fait que personne
n'aurait énoncé.

Ces limites ne sont pas des défauts à corriger mais des propriétés du dispositif. Les
énoncer avant la présentation des résultats est la condition d'une lecture avertie de
ce qui suit.